% !PRINTED 313
\weberpage{313}

Wir stellen für die Anwendung die sieben Tripel zusammen, aus denen man die Quadrupel als die Ergänzung zu dem vollen Systeme ablesen kann. In der ersten Columne steht das Zeichen ($\alpha$), zu dem das Tripel gehört:

\[
\begin{array}{c||ccc|}
1\,0\,0 & 0\,1\,0 & 0\,0\,1 & 0\,1\,1 \\
0\,1\,0 & 1\,0\,0 & 0\,0\,1 & 1\,0\,1 \\
0\,0\,1 & 1\,0\,0 & 0\,1\,0 & 1\,1\,0 \\
0\,1\,1 & 1\,0\,0 & 0\,1\,1 & 1\,1\,1 \\
1\,0\,1 & 0\,1\,0 & 1\,0\,1 & 1\,1\,1 \\
1\,1\,0 & 0\,0\,1 & 1\,1\,0 & 1\,1\,1 \\
1\,1\,1 & 0\,1\,1 & 1\,0\,1 & 1\,1\,0 \\
\end{array}
\]

Bezeichnen wir die Symbole ($\alpha$) in der Reihenfolge, in der sie in der ersten Columne dieser Tabelle stehen, mit $1, 2, 3, 4, 5, 6, 7$, so können wir die Tabelle für die Tripel und Quadrupel einfacher so darstellen:

\[
\begin{array}{c|ccc|cccc}
1 & 2 & 3 & 4 & 1 & 5 & 6 & 7 \\
2 & 1 & 3 & 5 & 2 & 4 & 6 & 7 \\
3 & 1 & 2 & 6 & 3 & 4 & 5 & 7 \\
4 & 1 & 4 & 7 & 2 & 3 & 5 & 6 \\
5 & 2 & 5 & 7 & 1 & 3 & 4 & 6 \\
6 & 3 & 6 & 7 & 1 & 2 & 4 & 5 \\
7 & 4 & 5 & 6 & 1 & 2 & 3 & 7 \\
\end{array}
\]

\begin{center}
\S.\,82.

\gesperrt{Anwendung auf Gleichungen $8^{\mathrm{ten}}$ Grades.}
\end{center}

In den abgeleiteten Formeln setzen wir nun für die Variablen $X_{z_1,z_2,z_3}$ die Wurzeln einer Gleichung $8^{\mathrm{ten}}$ Grades, von der wir annehmen, dass ihre \gesperrt{Galois}'sche Gruppe $P = TQ$ in dem festgesetzten Rationalitätsbereiche $\Omega$ in der linearen Congruenzgruppe $R = SQ$ enthalten sei, so dass $T$ ein Theiler von $S$ ist. Alle Functionen, die die Permutationen dieser Gruppe gestatten, sind rational.

Dazu gehört zunächst die Summe der $X$:

\weq{1}{\sum^{z} X_{z_1, z_2, z_3} = B,}
%CONT>>
